% GENERATED FILE. Edit canonical lessons, then run npm run generate:books.
% canonical-source-hash: fnv1a64:e575fbfcce35139f
% canonical-lessons: SA-C237-kumbhakarah, SA-C237-svarnakarah, SA-C237-lohakarah, SA-C237-tantuvayah, SA-C237-rajakah

\chapter{Potter, Goldsmith, Blacksmith, Weaver, Washerman}
\label{ch:sa-potter-goldsmith-blacksmith-weaver-washerman}
\hlchaptermodality{237}

\begin{chapteropening}
\textbf{By the end of this chapter:} \emph{I can say a potter, a goldsmith, a blacksmith, a weaver and a washerman.}

Complete the last lesson of chapter 237: I can say a potter, a goldsmith, a blacksmith, a weaver and a washerman.
\end{chapteropening}

\section[\texorpdfstring{\sk{कुम्भकारः} (kumbhakāraḥ)}{कुम्भकारः (kumbhakāraḥ)}]{\sk{कुम्भकारः} (kumbhakāraḥ) — a potter}

\label{lesson:SA-C237-kumbhakarah}



\begin{quote}
Before the new one: say the Sanskrit for a young woman, then the Sanskrit for a person.
\end{quote}

\subsection*{You'll want to know: \sk{कुम्भकारः}}
\textbf{\sk{कुम्भकारः}} — \emph{kumbhakāraḥ} — \textquotedblleft{}a potter\textquotedblright{}.

\begin{grammarlens}[title={What you've built}]
One more word for this chapter.
\end{grammarlens}

\subsection*{Your turn}
\begin{itemize}
\raggedright
  \item \emph{Say it:} \emph{kumbhakāraḥ}
  \item \emph{Say it:} \emph{kumbhakāraḥ}, once more
  \item \emph{Say it:} say \emph{kumbhakāraḥ}
  \item \emph{Recall:} say the Sanskrit for a young woman, then the Sanskrit for a person, then say \emph{kumbhakāraḥ} again
\end{itemize}

\begin{tcolorbox}[breakable,colback=teal!4,colframe=teal!35!black,title={Before you move on}]
What does \sk{कुम्भकारः} mean? (\textquotedblleft{}A potter\textquotedblright{}.) Say it once more.
\end{tcolorbox}
\section[\texorpdfstring{\sk{स्वर्णकारः} (svarṇakāraḥ)}{स्वर्णकारः (svarṇakāraḥ)}]{\sk{स्वर्णकारः} (svarṇakāraḥ) — a goldsmith}

\label{lesson:SA-C237-svarnakarah}



\begin{quote}
Before the new one: say the Sanskrit for a person, then the Sanskrit for a potter.
\end{quote}

\subsection*{You'll want to know: \sk{स्वर्णकारः}}
\textbf{\sk{स्वर्णकारः}} — \emph{svarṇakāraḥ} — \textquotedblleft{}a goldsmith\textquotedblright{}.

\begin{grammarlens}[title={What you've built}]
One more word for this chapter.
\end{grammarlens}

\subsection*{Your turn}
\begin{itemize}
\raggedright
  \item \emph{Say it:} \emph{svarṇakāraḥ}
  \item \emph{Say it:} \emph{svarṇakāraḥ}, once more
  \item \emph{Say it:} say \emph{svarṇakāraḥ}
  \item \emph{Recall:} say the Sanskrit for a person, then the Sanskrit for a potter, then say \emph{svarṇakāraḥ} again
\end{itemize}

\begin{tcolorbox}[breakable,colback=teal!4,colframe=teal!35!black,title={Before you move on}]
What does \sk{स्वर्णकारः} mean? (\textquotedblleft{}A goldsmith\textquotedblright{}.) Say it once more.
\end{tcolorbox}
\section[\texorpdfstring{\sk{लोहकारः} (lohakāraḥ)}{लोहकारः (lohakāraḥ)}]{\sk{लोहकारः} (lohakāraḥ) — a blacksmith}

\label{lesson:SA-C237-lohakarah}



\begin{quote}
Before the new one: say the Sanskrit for a potter, then the Sanskrit for a goldsmith.
\end{quote}

\subsection*{You'll want to know: \sk{लोहकारः}}
\textbf{\sk{लोहकारः}} — \emph{lohakāraḥ} — \textquotedblleft{}a blacksmith\textquotedblright{}.

\begin{grammarlens}[title={What you've built}]
One more word for this chapter.
\end{grammarlens}

\subsection*{Your turn}
\begin{itemize}
\raggedright
  \item \emph{Say it:} \emph{lohakāraḥ}
  \item \emph{Say it:} \emph{lohakāraḥ}, once more
  \item \emph{Say it:} say \emph{lohakāraḥ}
  \item \emph{Recall:} say the Sanskrit for a potter, then the Sanskrit for a goldsmith, then say \emph{lohakāraḥ} again
\end{itemize}

\begin{tcolorbox}[breakable,colback=teal!4,colframe=teal!35!black,title={Before you move on}]
What does \sk{लोहकारः} mean? (\textquotedblleft{}A blacksmith\textquotedblright{}.) Say it once more.
\end{tcolorbox}
\section[\texorpdfstring{\sk{तन्तुवायः} (tantuvāyaḥ)}{तन्तुवायः (tantuvāyaḥ)}]{\sk{तन्तुवायः} (tantuvāyaḥ) — a weaver}

\label{lesson:SA-C237-tantuvayah}



\begin{quote}
Before the new one: say the Sanskrit for a goldsmith, then the Sanskrit for a blacksmith.
\end{quote}

\subsection*{You'll want to know: \sk{तन्तुवायः}}
\textbf{\sk{तन्तुवायः}} — \emph{tantuvāyaḥ} — \textquotedblleft{}a weaver\textquotedblright{}.

\begin{grammarlens}[title={What you've built}]
One more word for this chapter.
\end{grammarlens}

\subsection*{Your turn}
\begin{itemize}
\raggedright
  \item \emph{Say it:} \emph{tantuvāyaḥ}
  \item \emph{Say it:} \emph{tantuvāyaḥ}, once more
  \item \emph{Say it:} say \emph{tantuvāyaḥ}
  \item \emph{Recall:} say the Sanskrit for a goldsmith, then the Sanskrit for a blacksmith, then say \emph{tantuvāyaḥ} again
\end{itemize}

\begin{tcolorbox}[breakable,colback=teal!4,colframe=teal!35!black,title={Before you move on}]
What does \sk{तन्तुवायः} mean? (\textquotedblleft{}A weaver\textquotedblright{}.) Say it once more.
\end{tcolorbox}
\section[\texorpdfstring{\sk{रजकः} (rajakaḥ)}{रजकः (rajakaḥ)}]{\sk{रजकः} (rajakaḥ) — a washerman}

\label{lesson:SA-C237-rajakah}



\begin{quote}
Before the new one: say the Sanskrit for a blacksmith, then the Sanskrit for a weaver.
\end{quote}

\subsection*{You'll want to know: \sk{रजकः}}
\textbf{\sk{रजकः}} — \emph{rajakaḥ} — \textquotedblleft{}a washerman\textquotedblright{}.

\begin{grammarlens}[title={What you've built}]
One more word for this chapter.
\end{grammarlens}

\subsection*{Your turn}
\begin{itemize}
\raggedright
  \item \emph{Say it:} \emph{rajakaḥ}
  \item \emph{Say it:} \emph{rajakaḥ}, once more
  \item \emph{Say it:} say \emph{rajakaḥ}
  \item \emph{Recall:} say the Sanskrit for a blacksmith, then the Sanskrit for a weaver, then say \emph{rajakaḥ} again
\end{itemize}

\begin{tcolorbox}[breakable,colback=teal!4,colframe=teal!35!black,title={Before you move on}]
What does \sk{रजकः} mean? (\textquotedblleft{}A washerman\textquotedblright{}.) Say it once more.
\end{tcolorbox}
